\subsection{The modules of a lax localization}
\label{subsec:laxloc-description}

\gutentag{00DB}%
Let $L \dashv R$ be a lax localization. By
\cref{thm:kz-modules,thm:kz-module-maps}, the $T$-modules form the locally full
sub-$(\infty,2)$-category
$\D_{T} \subset \D$
whose objects are the $x$ with $\eta_{x}$ reflective and whose $1$-morphisms
are those with left Beck--Chevalley unit square, and
$\iota_{1}\D_{T} \simeq \modcat{T}{\D}$.
For $x \in \D_{T}$ we write
$a_{x}\colon Tx \to x$
for the left adjoint of $\eta_{x}$, the action.
Let
$\Sigma_{\eta} \subset \func(\Delta^{1},\D)$
be the locally full sub-$(\infty,2)$-category generated by the $\eta_{d}$ and
their naturality squares
$(g,Tg)\colon \eta_{d} \to \eta_{d'}$
for $g\colon d \to d'$.

\begin{defi}
\label{def:closed-loc}
\gutentag{00DC}%
A locally full sub-$(\infty,2)$-category
$\B \subset \D$
is \emph{closed under localizations} if:
\begin{enumerate}
\item[(R1)] if $d$ is in $\B$ and $i\colon e \to d$ is reflective, then $e$ and
$i^{L}$ are in $\B$;
\item[(R2)] if $i\colon e \to d$ is reflective and $i^{L}$ and $fi^{L}$ are in
$\B$, then $f$ is in $\B$.
\end{enumerate}
\end{defi}

\gutentag{00DD}%
By \cref{lem:laxloc-closure}, every $\Pc_{\emptyset,P}$ is closed under
localizations.

\begin{thm}
\label{thm:laxloc-local}
\gutentag{00DE}%
Let $L \dashv R$ be a lax localization.
\begin{enumerate}
\item $R$ factors through $\D_{T}$.
\item
$\D_{T} = \Pc_{\emptyset,\Sigma_{\eta}} =
\Pc_{\emptyset,\{\eta_{d}\}_{d \in \D}}$.
\item $\D_{T}$ is the smallest locally full sub-$(\infty,2)$-category of $\D$
that contains the free modules $Td$ and the $1$-morphisms $Tg$ and is closed
under localizations. In particular, an object of $\D$ lies in $\D_{T}$ if and
only if
it admits
a reflective $1$-morphism to some $Rc$.
\end{enumerate}
\end{thm}

\begin{proof}
\gutentag{00DF}%
(1) By \cref{lem:laxloc-equiv}(3),
$R\varepsilon \dashv \eta R$
is a reflection in $\func(\C,\D)$. By \cref{lem:mates-refl}(2), every
$\eta_{Rc}$ is reflective and the naturality squares of $\eta R$, which are the
unit squares of the $Rg$, are left Beck--Chevalley. So $R$ factors through
$\D_{T}$ \cite[Cor.~2.6.6]{abellangagnahaugseng2024straightening}.

(2) By \cref{lem:laxloc-equiv}(2) and \cref{lem:mates-refl}(2), every
$L\eta_{d}$ is coreflective and $L$ takes the squares $(g,Tg)$ to right
Beck--Chevalley squares, so every $Rc$ and $Rg$ is $\Sigma_{\eta}$-local by
\cref{prop:laxloc-adjunction}. We show
$\D_{T} \subset \Pc_{\emptyset,\Sigma_{\eta}} \subset
\Pc_{\emptyset,\{\eta_{d}\}} \subset \D_{T}$,
where the middle inclusion is clear. For $x \in \D_{T}$,
\cref{lem:laxloc-closure}(2)
for the reflective
$\eta_{x}\colon x \to Tx = RLx$
shows that $x$ and $a_{x}$ are $\Sigma_{\eta}$-local. For
$f\colon x \to y$
in $\D_{T}$, the left mate of its unit square is an equivalence
$a_{y} \circ Tf \simeq f \circ a_{x}$,
so $fa_{x}$ is $\Sigma_{\eta}$-local, and so is $f$ by
\cref{lem:laxloc-closure}(3).

For the last inclusion, let $\Lambda^{z}_{d}$ denote the left adjoint of
$\eta_{d}^{\natural}\colon \Hom_{\D}(Td,z) \to \Hom_{\D}(d,z)$
for $\{\eta_{d}\}$-local $z$. For $z = Ty$, \eqref{eq:laxloc-adjunction}
identifies $\eta_{d}^{\natural}$ with
$(L\eta_{d})^{\natural}\colon \Hom_{\C}(LTd,Ly) \to \Hom_{\C}(Ld,Ly)$,
whose left adjoint is
$\varepsilon_{Ld}^{\natural}$
by
\cref{lem:laxloc-equiv}(4). As
$m \circ \varepsilon_{Ld}$
corresponds to
$Rm \circ \mu_{d} \circ \eta_{Td} \simeq Rm$,
the essential image of $\Lambda^{Ty}_{d}$ consists of the $Rm$ for
$m\colon Ld \to Ly$.

Let $x$ be $\{\eta_{d}\}$-local, let
$a := \Lambda^{x}_{x}(\mathrm{id}_{x})$,
and let
$c\colon a\eta_{x} \simeq \mathrm{id}_{x}$
be the inverse of the unit. As
$\mathrm{id}_{Tx} = R(\mathrm{id}_{Lx})$,
we have
$\mathrm{id}_{Tx} \simeq \Lambda^{Tx}_{x}(\eta_{x})$,
so the $2$-morphism
$\eta_{x}c^{-1}\colon \eta_{x} \Rightarrow \eta_{x}^{\natural}(\eta_{x}a)$
has an adjunct
$\nu\colon \mathrm{id}_{Tx} \Rightarrow \eta_{x}a$.
Then $\nu\eta_{x}$ is invertible, as the unit of
$\Lambda^{Tx}_{x} \dashv \eta_{x}^{\natural}$
is. Both $a$ and
$a\eta_{x}a \simeq a$
lie in the essential image of $\Lambda^{x}_{x}$, on which $\eta_{x}^{\natural}$
is fully faithful, and
$\eta_{x}^{\natural}(a\nu) = a\nu\eta_{x}$
is invertible; so $a\nu$ is invertible, and
$a \dashv \eta_{x}$
with counit $c$ by \cref{lem:laxloc-criterion}(1). So $x \in \D_{T}$ and
$a_{x} = a$.

Let $f\colon x \to y$ be $\{\eta_{d}\}$-local. The left adjoints $a_{x}$ and
$a_{y}$ are $\{\eta_{d}\}$-local by \cref{lem:laxloc-closure}(1). As
$a_{x} = \Lambda^{x}_{x}(\mathrm{id}_{x})$
and $Tf = R(Lf)$ lies in the essential image of $\Lambda^{Ty}_{x}$,
the criterion
after \cref{def:laxloc-local} shows that $f \circ a_{x}$ and
$a_{y} \circ Tf$
lie in the essential image of $\Lambda^{y}_{x}$. So the left mate
$a_{y} \circ Tf \Rightarrow f \circ a_{x}$
of the unit square of $f$ is invertible if its whiskering with $\eta_{x}$ is.
By \cref{constr:mates}, this whiskering is
$a_{y}Tf\,u_{x}\eta_{x}$
composed
with invertible $2$-morphisms, where $u_{x}$ is the unit of
$a_{x} \dashv \eta_{x}$,
and $u_{x}\eta_{x}$ is invertible by a triangle identity. So $f$ is in
$\D_{T}$.

(3) By (1) and (2), $\D_{T}$ contains the $Rc$ and $Rg$, in particular the $Td$
and $Tg$, and is closed under localizations. Conversely, let $\B$ be such. For
$x \in \D_{T}$, (R1) for
$\eta_{x}\colon x \to Tx$
shows that $x$ and $a_{x}$ are in $\B$, and for $f\colon x \to y$ in $\D_{T}$,
(R2) for
$fa_{x} \simeq a_{y} \circ Tf$,
which is in $\B$, shows that $f$ is in $\B$. The last claim follows from (R1)
for $\D_{T}$, as
$\eta_{x}\colon x \to RLx$
is reflective for $x \in \D_{T}$.
\end{proof}

\begin{rmk}
\label{rmk:laxloc-local}
\gutentag{00E0}%
For $2$-categories, the equality
$\D_{T} = \Pc_{\emptyset,\{\eta_{d}\}}$
is \cite[Thm.~2.7]{dilibertilobbiasousa2022kz}, based on Marmolejo--Wood
\cite{marmolejowood2012kan}, and \cref{lem:laxloc-closure}(2), (3) are
contained in \cite[Lemma~2.6]{dilibertilobbiasousa2022kz}. By (3), $\D_{T}$ is
the closure of the free modules under localizations, just as the local objects
of a Bousfield localization are the closure of the free modules under
equivalences.
\end{rmk}

\begin{prop}
\label{prop:laxloc-normalize}
\gutentag{00E1}%
Let $L \dashv R$ be a lax localization, and
$\iota\colon \D_{T} \to \D$
the inclusion. Then
$T \simeq \iota L_{T}$
for a functor
$L_{T}\colon \D \to \D_{T}$,
and
$L_{T} \dashv \iota$
with unit $\eta$, and its counit has components $a_{x}$. This adjunction is a
lax localization, whose monad has underlying functor $T$ and unit $\eta$.
\end{prop}

\begin{proof}
\gutentag{00E2}%
By \cref{thm:laxloc-local}(1), $T = RL$ factors as $\iota L_{T}$. The
components $\eta_{x}$ of
$\eta\iota\colon \iota \Rightarrow T\iota$
are reflective and its naturality squares, the unit squares of the
$1$-morphisms of $\D_{T}$, are left Beck--Chevalley. So by
\cite[Cor.~5.2.12]{abellangagnahaugseng2024straightening} (\cref{lem:mates}(2))
it has a left adjoint $a$ in $\func(\D_{T},\D)$, with components $a_{x}$ and
invertible counit. The $a_{x}$ are in $\D_{T}$ by \cref{lem:laxloc-closure}(1)
and \cref{thm:laxloc-local}(2), so $a$ lifts to
$L_{T}\iota \Rightarrow \mathrm{id}_{\D_{T}}$,
as $\func(\D_{T},-)$ preserves locally full inclusions
(\cref{subsec:prelim-conventions}). For the triangle identities,
$a \circ \eta\iota \simeq \mathrm{id}_{\iota}$
is the counit of
$a \dashv \eta\iota$,
and
$\iota(aL_{T} \circ L_{T}\eta) \simeq aL_{T} \circ T\eta \simeq \mu \circ T\eta
\simeq \mathrm{id}_{T}$,
as $aL_{T}$ and $\mu$ are both left adjoint to $\eta T$
(\cref{lem:kz-basic}(1)); this lifts along the locally fully faithful
$\iota_{\natural}\colon \func(\D,\D_{T}) \to \func(\D,\D)$.
So $\eta$ and the lift of $a$ are the unit and counit of an adjunction
$L_{T} \dashv \iota$
(\cref{subsec:prelim-adjunctions}).
Finally, \cref{lem:laxloc-equiv}(5) holds for
$L_{T} \dashv \iota$,
as
$a_{x} \dashv \eta_{x}$
with counit the triangle identity.
\end{proof}

\begin{defi}
\label{def:laxloc-normalized}
\gutentag{00E3}%
A lax localization is \emph{normalized} if $R$ factors through an equivalence
$\C \simeq \D_{T}$.
\end{defi}

\gutentag{00E4}%
By \cref{thm:laxloc-local}, a lax localization is normalized if and only if $R$
is a locally full inclusion whose image is closed under localizations, and
\cref{prop:laxloc-normalize} replaces every lax localization by a normalized one
with the same monad functor and unit. Normalization is not automatic:

\begin{exmp}
\label{ex:laxloc-initial}
\gutentag{00E5}%
If $\C$ has an object $0$ with
$\Hom_{\C}(0,c) \simeq \mathrm{pt}$
for all $c$, then
$0\colon \mathrm{pt} \to \C$
is left adjoint to
$\C \to \mathrm{pt}$,
with $T = \mathrm{id}$. This is a lax localization with
$\D_{T} = \mathrm{pt}$,
normalized only if
$\C \simeq \mathrm{pt}$.
\end{exmp}

\gutentag{00E6}%
Let
$\Sat_{L} \subset \func(\Delta^{1},\D)$
be the preimage under $L$ of the image of
$\ell^{*}\colon \func(\Corefl,\C) \to \func(\Delta^{1},\C)$:
its objects are the $1$-morphisms $f$ with $Lf$ coreflective, and its
$1$-morphisms are the commutative squares $\sigma$ between them,
drawn with these
vertical, with $L\sigma$ right Beck--Chevalley (\cref{lem:mates-refl}). It is
locally full, as locally full inclusions are stable under pullback.

\begin{prop}
\label{prop:laxloc-class}
\gutentag{00E7}%
Let $L \dashv R$ be a lax localization.
\begin{enumerate}
\item
$\Sigma_{\eta} \subset \Sat_{L}$.
\item A $1$-morphism $f$ is in $\Sat_{L}$ if and only if
$\D_{T} \subset \Pc_{\emptyset,\{f\}}$,
and a commutative square $\sigma$ between $1$-morphisms in $\Sat_{L}$ is an
$\Sat_{L}$-square if and only if
$\D_{T} \subset \Pc_{\emptyset,\{\sigma\}}$.
\item
$\Pc_{\emptyset,\Sat_{L}} = \D_{T}$,
and $\Sat_{L}$ is the largest locally full sub-$(\infty,2)$-category $P$ of
$\func(\Delta^{1},\D)$ with
$\D_{T} \subset \Pc_{\emptyset,P}$.
\end{enumerate}
\end{prop}

\begin{proof}
\gutentag{00E8}%
(1) was shown in the proof of \cref{thm:laxloc-local}(2). (2) By
\cref{prop:laxloc-adjunction}, $f$ is in $\Sat_{L}$ if and only if every $Rc$
and $Rg$ lies in
$\Pc_{\emptyset,\{f\}}$,
and as
$\Pc_{\emptyset,\{f\}}$
is closed under localizations, this holds if and only if
$\D_{T} \subset \Pc_{\emptyset,\{f\}}$
by \cref{thm:laxloc-local}(1), (3). Likewise for $\sigma$.
(3) By (2),
$\D_{T} \subset \Pc_{\emptyset,\Sat_{L}}$,
and by (1) and \cref{thm:laxloc-local}(2),
$\Pc_{\emptyset,\Sat_{L}} \subset \Pc_{\emptyset,\Sigma_{\eta}} = \D_{T}$.
If
$\D_{T} \subset \Pc_{\emptyset,P}$,
then
$\D_{T} \subset \Pc_{\emptyset,\{\sigma\}}$
for every object and $1$-morphism $\sigma$ of $P$, so
$P \subset \Sat_{L}$
by (2).
\end{proof}

\begin{rmk}
\label{rmk:laxloc-saturated}
\gutentag{00E9}%
By \cref{prop:laxloc-class}(2) and \cref{lem:saturation-explicit}, which we will
prove in \cref{subsec:2dloc-saturation},
$\Sat_{L} = \Sat^{\mathrm{corefl}}_{\D_{T}}$,
so $\Sat_{L}$ is coreflectively saturated by \cref{thm:saturation-saturated},
without any presentability; and
$\Sat_{L} = \overline{\Sigma_{\eta}} = \overline{P}$
for every $P$ with
$\Pc_{\emptyset,P} = \D_{T}$
(\cref{def:saturation}). For an arbitrary adjunction,
\cref{prop:laxloc-adjunction}
shows in the same way that $\Sat_{L}$ is the coreflective saturation of the
locally full sub-$(\infty,2)$-category generated by the image of $R$, so it is
coreflectively saturated as well.
\end{rmk}
